\documentclass[11pt,twoside]{article}

\usepackage[
  paperwidth=8.5in, paperheight=11in,
  twoside,
  inner=0.95in, outer=2.85in,
  top=1.0in, bottom=1.0in,
  marginparwidth=2.2in, marginparsep=0.2in,
  headsep=16pt
]{geometry}

\usepackage{annotated}
\usepackage[colorlinks=true, allcolors=RuInk]{hyperref}

\linespread{1.05}
\setlength{\parskip}{0pt}
\setlength{\parindent}{1.4em}

\renewcommand{\leftmark}{Chapter 1 \;\textbullet\; The Real and Complex Number Systems}
\renewcommand{\rightmark}{The Annotated Rudin}

\begin{document}

% ============================================================
\thispagestyle{empty}
\vspace*{1.5in}
\begin{center}
{\sffamily\Huge\bfseries\color{RuInk} The Annotated Rudin}\\[10pt]
{\sffamily\large\color{RuGrey} \emph{Principles of Mathematical Analysis}, Third Edition}\\[4pt]
{\sffamily\large\color{RuGrey} Walter Rudin, with a running commentary}\\[40pt]
{\sffamily\normalsize\bfseries\color{RuInk} STYLE SAMPLE \textperiodcentered\ ADDITIVE \textperiodcentered\ B/W PRINT}\\[6pt]
{\sffamily\small\color{RuGrey} Chapter 1, \S\S1.1--1.11 and 1.19--1.21}
\end{center}
\vfill
\begin{center}
\begin{minipage}{0.8\textwidth}
\small\color{RuGrey}
Rudin's text appears here complete and unaltered --- his sentences, his
paragraph breaks, his numbering. Nothing has been paraphrased away. The
annotation is strictly additive: numbered markers in his prose key to notes in
the outer margin, and longer commentary sits in ruled boxes between his
paragraphs, where it can be skipped without breaking his argument.
\end{minipage}
\end{center}
\newpage

% ============================================================
\thispagestyle{empty}
\section*{How to read this book}

Everything set in the main column below is Rudin, verbatim. You can read it
straight through and never look at an annotation. The annotation layer sits in
two places, both outside his text.

\medskip
\noindent
\textbf{Margin notes.} A small superscript number in Rudin's text --- like
this\kn{The marker is the only mark added to his page, and it is the same for
every tier so his line spacing is undisturbed. The number out here tells you
what kind of note it is.} --- keys to a numbered note in the outer margin. Since
the book prints in black and white, the \emph{tier} of a note is carried by how
its number is set:

\medskip
\begin{tabular}{@{}p{0.15\textwidth}p{0.78\textwidth}@{}}
\tierplain{7} & \textbf{Plain number --- unpacking.} What the sentence is
actually saying, where a formula came from, an omitted algebraic step.\\[5pt]
\tierdeep{7} & \textbf{Outlined number --- going deeper.} Context beyond the
course: what breaks in general, which later theorem this sets up. Skippable.\\[5pt]
\tierwarn{7} & \textbf{Reversed number --- warning.} A specific belief students
form here that is false, or a place where Rudin's wording will mislead you.
\end{tabular}

\medskip
\noindent
\textbf{Interleaved boxes.} When a note runs longer than a margin will hold, it
becomes a box between Rudin's paragraphs, framed by the same three tiers: a
\textbf{solid} left rule for unpacking, a \textbf{dashed} left rule for going
deeper, and rules \textbf{top and bottom} with a black tab for a warning.

\medskip
\noindent
\textbf{Structural commentary.} Two further boxes carry no mathematics at all.
They explain the book rather than its contents --- why a chapter exists, why
Rudin ordered it as he did, why a section that looks like a detour is not one.

\begin{roadmap}[Heavy frame, black title bar]
Opens and closes every chapter. Says what the chapter is for, how it is
organized and why in that order, and which parts are load-bearing versus
bookkeeping. Where Rudin has explained his own choices --- usually in the
Preface --- these boxes quote him rather than guess.
\end{roadmap}

\begin{signpost}[Thin frame, inline heading]
Opens a section. Shorter: why this material sits at this point, what Rudin is
setting up, and what to hold on to while reading it.
\end{signpost}

\newpage

% ============================================================
\begin{roadmap}[Why this chapter exists]
Analysis is the study of limits, and a limit is a promise that some process
arrives somewhere. The rationals break that promise: you can write down a
sequence of rationals closing in on $\sqrt2$ with nothing at the end of it.
Every theorem in this book --- convergence, continuity, the intermediate value
theorem, the existence of integrals --- would be false or vacuous over $\Q$.

\medskip
So before any analysis can begin, one fact must be secured: \emph{there is a
number system with no gaps}. Chapter 1 secures it, and does nothing else. Its
entire output is Theorem 1.19, plus enough consequences (1.20, 1.21) to show the
theorem was worth having.
\end{roadmap}

\begin{roadmap}[How it is organized, and why in that order]
Rudin's arrangement is deliberate, and he explains it in the Preface. The
obvious order would be to construct $\R$ from $\Q$ first and then do analysis on
it. He refuses, on the grounds that beginning with the construction is
``pedagogically unsound (though logically correct)'' because ``most students
simply fail to appreciate the need for doing this.'' The chapter is therefore
built to make you \emph{feel the need before paying the cost}:

\medskip
\centering
\begin{tikzpicture}[scale=0.92, transform shape, node distance=5mm]
\node[rmbox, text width=22mm] (a) {\textbf{1.1--1.2}\\exhibit the gap\\in $\Q$};
\node[rmbox, text width=22mm, right=of a] (b) {\textbf{1.5--1.18}\\vocabulary:\\order, field};
\node[rmbox, text width=22mm, right=of b] (c) {\textbf{1.19}\\assert $\R$\\(proof deferred)};
\node[rmbox, text width=22mm, right=of c] (d) {\textbf{1.20--1.21}\\cash it out\\at once};
\draw[rmarrow] (a) -- (b); \draw[rmarrow] (b) -- (c); \draw[rmarrow] (c) -- (d);
\end{tikzpicture}

\medskip
\raggedright
Three consequences of this choice are worth knowing in advance.

\smallskip
\textbf{The construction is missing on purpose.} Theorem 1.19 is stated and not
proved; the Dedekind-cut construction sits in the appendix, to be read ``whenever
the time seems ripe.'' Nothing later depends on having read it.

\smallskip
\textbf{The abstraction comes before the object.} Rudin defines ordered sets
(1.5--1.11) and fields (1.12--1.18) \emph{in general} before naming $\R$. This
looks like a detour and is not: it lets him characterise $\R$ by properties
rather than by construction, and it means $\C$ and $\R^k$ inherit the same
axioms later without a second round of proofs. He says so himself at 1.13(c) ---
``once we do this, we will not need to do it again.''

\smallskip
\textbf{The tail of the chapter is bookkeeping.} Everything after 1.21 --- the
extended reals, $\C$, Euclidean spaces, Schwarz's inequality --- is notation
being laid in for later chapters, not argument. Read it once and move on.
\end{roadmap}

\newpage

% ============================================================
\begin{center}
{\sffamily\LARGE\bfseries\color{RuInk} Chapter 1}\\[6pt]
{\sffamily\Large\bfseries\color{RuInk} THE REAL AND COMPLEX NUMBER SYSTEMS}
\end{center}

\rsection{Introduction}

\noindent
A satisfactory discussion of the main concepts of analysis (such as convergence,
continuity, differentiation, and integration) must be based on an accurately
defined number concept. We shall not, however, enter into any discussion of the
axioms that govern the arithmetic of the integers, but assume familiarity with
the rational numbers (i.e., the numbers of the form $m/n$ where $m$ and $n$ are
integers and $n \ne 0$).\kn{Rudin is choosing a starting line, and he defends
the choice in the Preface: beginning with a construction of $\R$ is
``pedagogically unsound (though logically correct)'' because ``most students
simply fail to appreciate the need for doing this.'' So the construction is
exiled to the appendix, and $\R$ will arrive in 1.19 as an assertion. This
sentence marks everything he is willing to take for granted.}

The rational number system is inadequate for many purposes, both as a field and
as an ordered set. (These terms will be defined in Definitions 1.6 and 1.12.)
For instance, there is no rational $p$ such that $p^2 = 2$. (We shall prove this
presently.) This leads to the introduction of so-called ``irrational numbers''
which are often written as infinite decimal expansions and are considered to be
``approximated'' by the corresponding finite decimals. Thus the sequence
\[ 1,\ 1.4,\ 1.41,\ 1.414,\ 1.4142,\ \dots \]
``tends to $\sqrt 2$.''\kw{The scare quotes are doing real work. ``Tends to'' is
undefined until Definition 3.1, and the thing it tends to is undefined until
Theorem 1.19. Rudin is pointing at a circularity he is about to break, not
relying on one.} But unless the irrational number $\sqrt 2$ has been clearly
defined, the question must arise: Just what is it that this sequence ``tends
to''?

This sort of question can be answered as soon as the so-called ``real number
system'' is constructed.

\ritem{1.1}{Example}
We now show that the equation
\begin{equation}
  p^2 = 2 \tag{1.1}
\end{equation}
is not satisfied by any rational $p$. If there were such a $p$, we could write
$p = m/n$ where $m$ and $n$ are integers that are not both even.\kn{Why we may
assume this: put $p$ in lowest terms. If $m$ and $n$ were both even we could
cancel a factor of $2$ and repeat, and the process must stop because $|m|$
strictly decreases. That appeal to ``must stop'' is well-ordering, so the step
is not quite free.} Let us assume this is done. Then (1.1) implies
\begin{equation}
  m^2 = 2n^2, \tag{1.2}
\end{equation}
This shows that $m^2$ is even. Hence $m$ is even (if $m$ were odd, $m^2$ would
be odd), and so $m^2$ is divisible by $4$. It follows that the right side of
(1.2) is divisible by $4$, so that $n^2$ is even, which implies that $n$ is
even.

The assumption that (1.1) holds thus leads to the conclusion that both $m$ and
$n$ are even, contrary to our choice of $m$ and $n$. Hence (1.1) is impossible
for rational $p$.

We now examine this situation a little more closely. Let $A$ be the set of all
positive rationals $p$ such that $p^2 < 2$ and let $B$ consist of all positive
rationals $p$ such that $p^2 > 2$. We shall show that \emph{$A$ contains no
largest number and $B$ contains no smallest}.\kn{This, and not the irrationality
of $\sqrt2$, is the real content of \S1.1. Irrationality says
$\sqrt2 \notin \Q$. This says $\Q$ has a \emph{gap} where $\sqrt 2$ should be
--- a structural defect, and the one Definition 1.10 will name.}

More explicitly, for every $p$ in $A$ we can find a rational $q$ in $A$ such
that $p < q$, and for every $p$ in $B$ we can find a rational $q$ in $B$ such
that $q < p$.

To do this, we associate with each rational $p > 0$ the number
\begin{equation}
  q = p - \frac{p^2-2}{p+2} = \frac{2p+2}{p+2}. \tag{1.3}
\end{equation}
Then
\begin{equation}
  q^2 - 2 = \frac{2(p^2-2)}{(p+2)^2}. \tag{1.4}
\end{equation}
If $p$ is in $A$, then $p^2 - 2 < 0$, (1.3) shows that $q > p$ and (1.4) shows
that $q^2 < 2$. Thus $q$ is in $A$.

If $p$ is in $B$, then $p^2 - 2 > 0$, (1.3) shows that $0 < q < p$, and (1.4)
shows that $q^2 > 2$. Thus $q$ is in $B$.

\begin{unpack}[Where (1.3) comes from]
Rudin produces this formula with no derivation, and it is the first place in the
book where a reader stalls. It is not Newton's method. Ask instead: which map
$p \mapsto q$ has $\sqrt2$ as a \emph{fixed point}? Try the simplest shape, a
ratio of linear terms $q = (ap+b)/(cp+d)$. Setting $q = p$ gives
$p(cp+d) = ap+b$, and we want that to collapse to $p^2 = 2$. The
smallest-integer solution is $q = (2p+2)/(p+2)$: it gives $p^2 + 2p = 2p+2$,
that is, $p^2 = 2$.

Why it moves $p$ toward $\sqrt2$ without crossing it becomes visible once you
factor the error:
\[
  q - \sqrt2 \;=\; \frac{2-\sqrt2}{p+2}\,\bigl(p - \sqrt2\bigr).
\]
For $p>0$ the multiplier lies strictly between $0$ and
$(2-\sqrt2)/2 \approx 0.29$. That it is \emph{positive} keeps $q$ on the same
side of $\sqrt2$ as $p$ --- Rudin's ``$q$ is in $A$'' and ``$q$ is in $B$.'' That
it is \emph{less than one} makes the error shrink, so $q$ is strictly better
than $p$ and no member of $A$ can be largest.

Equation (1.4) is worth checking rather than accepting:
$(2p+2)^2 - 2(p+2)^2 = (4p^2+8p+4)-(2p^2+8p+8) = 2p^2-4$.
\end{unpack}

\begin{deeper}[A contraction, eight chapters early]
That factorisation says $p \mapsto (2p+2)/(p+2)$ shrinks the distance to
$\sqrt2$ by a factor of at most $0.29$ every step --- exactly the hypothesis of
the contraction principle, Theorem 9.23, which Rudin later uses to prove the
inverse function theorem. Chapter 1 runs the argument by hand before the general
theorem exists.

Iterating from $p_0 = 1$ gives $1, \tfrac43, \tfrac{10}{7}, \tfrac{24}{17},
\tfrac{58}{41}, \dots$, the convergents of $\sqrt2 = [1;2,2,2,\dots]$. Newton's
method on $f(p) = p^2-2$ would give $p \mapsto (p^2+2)/(2p)$ instead: faster,
but it always lands \emph{above} $\sqrt2$, so it would prove ``$B$ has no
smallest'' and say nothing about $A$. Rudin's map does both halves at once.
\end{deeper}

\ritem{1.2}{Remark}
The purpose of the above discussion has been to show that the rational number
system has certain gaps, in spite of the fact that between any two rationals
there is another: if $r < s$, then $r < (r+s)/2 < s$.\kn{Density and
completeness are independent properties, and this sentence is the first warning
of it. $\Q$ is dense in itself yet riddled with holes; $\Z$ has no holes in the
relevant sense yet is nowhere dense. Neither property implies the other.} The
real number system fills these gaps. This is the principal reason for the
fundamental role which it plays in analysis.

In order to elucidate its structure, as well as that of the complex numbers, we
start with a brief discussion of the general concepts of \emph{ordered set} and
\emph{field}.

Here is some of the standard set-theoretic terminology that will be used
throughout this book.

\ritem{1.3}{Definitions}
If $A$ is any set (whose elements may be numbers or any other objects), we write
$x \in A$ to indicate that $x$ is a member (or an element) of $A$. If $x$ is not
a member of $A$, we write $x \notin A$.

The set which contains no element will be called the \emph{empty set}. If a set
has at least one element, it is called \emph{nonempty}.

If $A$ and $B$ are sets, and if every element of $A$ is an element of $B$, we
say that $A$ is a \emph{subset} of $B$, and we write $A \subset B$, or
$B \supset A$.\kw{Rudin's $\subset$ permits equality --- it is what most modern
texts write $\subseteq$. He writes ``proper subset'' when he means strict
containment. Read every $\subset$ in this book as $\subseteq$, or you will
misread hypotheses for the next 300 pages.} If, in addition, there is an element
of $B$ which is not in $A$, then $A$ is said to be a \emph{proper subset} of
$B$. Note that $A \subset A$ for every set $A$.

If $A \subset B$ and $B \subset A$, we write $A = B$. Otherwise, $A \ne B$.

\ritem{1.4}{Definition}
Throughout Chap.\ 1, the set of all rational numbers will be denoted by $\Q$.

\rsection{Ordered Sets}

\begin{signpost}[Why this section is here]
Nothing in this section mentions $\R$, and on a first reading it feels like an
unmotivated detour into abstraction. It is not. Rudin has just shown that $\Q$
fails in a specific way, and he now needs language precise enough to \emph{name}
that failure --- because the name will become the defining property of $\R$ in
1.19. Definitions 1.5 to 1.7 are only scaffolding; 1.8 and 1.10 are the payload.
If you read one thing here, read Definition 1.10, and notice that it is exactly
what Example 1.9(a) showed $\Q$ does not have.
\end{signpost}

\ritem{1.5}{Definition}
Let $S$ be a set. An \emph{order} on $S$ is a relation, denoted by $<$, with the
following two properties:
\begin{enumerate}[label=(\roman*), leftmargin=2.2em, itemsep=2pt, topsep=4pt]
  \item If $x \in S$ and $y \in S$, then one and only one of the statements
        \[ x < y, \qquad x = y, \qquad y < x \]
        is true.\kn{Trichotomy. ``One and only one'' packs in two separate
        demands: at least one holds, so the order is \emph{total}; and at most
        one holds, so $<$ is irreflexive and antisymmetric. Drop totality and
        you get a partial order, where suprema can fail to exist for reasons
        having nothing to do with gaps.}
  \item If $x, y, z \in S$, if $x < y$ and $y < z$, then $x < z$.
\end{enumerate}
The statement ``$x < y$'' may be read as ``$x$ is less than $y$'' or ``$x$ is
smaller than $y$'' or ``$x$ precedes $y$''.

It is often convenient to write $y > x$ in place of $x < y$.

The notation $x \le y$ indicates that $x < y$ or $x = y$, without specifying
which of these two is to hold. In other words, $x \le y$ is the negation of
$x > y$.\kd{That last sentence is a small theorem, not a definition: ``$x<y$ or
$x=y$'' is equivalent to ``not $x>y$'' precisely because of trichotomy in (i).
In a partial order the two come apart, which is why $\le$ is normally taken as
the primitive notion there.}

\ritem{1.6}{Definition}
An \emph{ordered set} is a set $S$ in which an order is defined.

For example, $\Q$ is an ordered set if $r < s$ is defined to mean that $s - r$
is a positive rational number.

\ritem{1.7}{Definition}
Suppose $S$ is an ordered set, and $E \subset S$. If there exists a
$\beta \in S$ such that $x \le \beta$ for every $x \in E$, we say that $E$ is
\emph{bounded above}, and call $\beta$ an \emph{upper bound} of $E$.\kn{$\beta$
is drawn from $S$, and need not lie in $E$. This is exactly why the next
definition has content: the upper bounds come from the ambient set, so a set can
be bounded above by elements sitting far away from it.}

Lower bounds are defined the same way (with $\ge$ in place of $\le$).

\ritem{1.8}{Definition}
Suppose $S$ is an ordered set, $E \subset S$, and $E$ is bounded above. Suppose
there exists an $\alpha \in S$ with the following properties:
\begin{enumerate}[label=(\roman*), leftmargin=2.2em, itemsep=2pt, topsep=4pt]
  \item $\alpha$ is an upper bound of $E$.
  \item If $\gamma < \alpha$ then $\gamma$ is not an upper bound of
        $E$.\kn{This is the clause you will actually use, and it is stated
        contrapositively. The working form: for every $\gamma < \alpha$ there
        is an $x \in E$ with $x > \gamma$. In $\R$, equivalently: for every
        $\varepsilon > 0$ there is $x \in E$ with $x > \alpha - \varepsilon$.
        Reach for that sentence every time a proof opens ``let
        $\alpha = \sup E$.''}
\end{enumerate}
Then $\alpha$ is called the \emph{least upper bound} of $E$ [that there is at
most one such $\alpha$ is clear from (ii)] or the \emph{supremum} of $E$, and we
write
\[ \alpha = \sup E. \]

The \emph{greatest lower bound}, or \emph{infimum}, of a set $E$ which is
bounded below is defined in the same manner: The statement
\[ \alpha = \inf E \]
means that $\alpha$ is a lower bound of $E$ and that no $\beta$ with
$\beta > \alpha$ is a lower bound of $E$.

\begin{unpack}[Why the least upper bound and not the largest element]
Because the largest element usually does not exist, and the supremum is the
repair. Take $E_1 = \{\, r \in \Q : r < 0 \,\}$. It has no largest element: for
any $r<0$ there is $r/2$ with $r < r/2 < 0$. But the set of its upper bounds ---
everything ${}\ge 0$ --- does have a smallest one, namely $0$.

The manoeuvre is to move the question from $E$, where the answer is often
``none,'' to the set of upper bounds of $E$, where it is often ``yes.'' Rudin's
bracketed aside is the payoff: uniqueness follows immediately from (ii), so
speaking of ``the'' supremum is legitimate.
\end{unpack}

\ritem{1.9}{Examples}
\begin{enumerate}[label=(\alph*), leftmargin=2.2em, itemsep=5pt, topsep=4pt]
  \item Consider the sets $A$ and $B$ of Example 1.1 as subsets of the ordered
        set $\Q$. The set $A$ is bounded above. In fact, the upper bounds of $A$
        are exactly the members of $B$.\kd{Not \emph{exactly}, and Rudin is
        being loose. Every $\beta \in B$ bounds $A$, and no member of $A$ does.
        But one must also rule out a positive rational with $\beta^2 = 2$, and
        that is impossible only by Example 1.1. The claim holds, but it leans on
        the irrationality result proved a page earlier.} Since $B$ contains no
        smallest member, $A$ has no least upper bound in $\Q$.

        Similarly, $B$ is bounded below: the set of all lower bounds of $B$
        consists of $A$ and of all $r \in \Q$ with $r \le 0$. Since $A$ has no
        largest member, $B$ has no greatest lower bound in $\Q$.
  \item If $\alpha = \sup E$ exists, then $\alpha$ may or may not be a member of
        $E$. For instance, let $E_1$ be the set of all $r \in \Q$ with $r < 0$.
        Let $E_2$ be the set of all $r \in \Q$ with $r \le 0$. Then
        \[ \sup E_1 = \sup E_2 = 0, \]
        and $0 \notin E_1$, $0 \in E_2$.\kn{Both suprema exist \emph{in $\Q$}.
        Part (a) fails not because $\Q$ is too small in general but because of
        where that particular gap sits. Whether $\sup E$ lies in $E$ is a
        separate question from whether $\sup E$ exists, and running the two
        together is the most common first-week error.}
  \item Let $E$ consist of all numbers $1/n$, where $n = 1,2,3,\dots$. Then
        $\sup E = 1$, which is in $E$, and $\inf E = 0$, which is not in $E$.
\end{enumerate}

\ritem{1.10}{Definition}
An ordered set $S$ is said to have the \emph{least-upper-bound property} if the
following is true:

If $E \subset S$, $E$ is not empty, and $E$ is bounded above, then $\sup E$
exists in $S$.\kw{Both hypotheses are load-bearing, and dropping either breaks
the definition rather than merely weakening it. In $\R$: every real is an upper
bound of $\emptyset$, so $\emptyset$ has no \emph{least} one; and $\Z$ is
unbounded above, so it has none either. Any proof that produces a supremum owes
you a check of both.}

Example 1.9(a) shows that $\Q$ does not have the least-upper-bound property.

We shall now show that there is a close relation between greatest lower bounds
and least upper bounds, and that every ordered set with the least-upper-bound
property also has the greatest-lower-bound property.

\ritem{1.11}{Theorem}
\emph{Suppose $S$ is an ordered set with the least-upper-bound property,
$B \subset S$, $B$ is not empty, and $B$ is bounded below. Let $L$ be the set of
all lower bounds of $B$. Then}
\[ \alpha = \sup L \]
\emph{exists in $S$ and $\alpha = \inf B$.}

\emph{In particular, $\inf B$ exists in $S$.}

\begin{proof}
Since $B$ is bounded below, $L$ is not empty. Since $L$ consists of exactly
those $y \in S$ which satisfy the inequality $y \le x$ for every $x \in B$, we
see that every $x \in B$ is an upper bound of $L$.\kn{This single sentence is
the whole idea. We are handed a hypothesis about suprema and asked for an
infimum, so we must find a set whose \emph{supremum} is the infimum we want.
$L$ is that set: $B$ presses down on it from above, and the two meet at one
point.} Thus $L$ is bounded above. Our hypothesis about $S$ implies therefore
that $L$ has a supremum in $S$; call it $\alpha$.

If $\gamma < \alpha$ then (see Definition 1.8) $\gamma$ is not an upper bound of
$L$, hence $\gamma \notin B$.\kn{Compressed to the point of opacity. Unpacked:
$\gamma$ failing to be an upper bound of $L$ produces some $y \in L$ with
$y > \gamma$. Since $y$ is a lower bound of $B$, every $x \in B$ satisfies
$x \ge y > \gamma$. So $\gamma$ itself cannot belong to $B$.} It follows that
$\alpha \le x$ for every $x \in B$. Thus $\alpha \in L$.

If $\alpha < \beta$ then $\beta \notin L$, since $\alpha$ is an upper bound of
$L$.

We have shown that $\alpha \in L$ but $\beta \notin L$ if $\beta > \alpha$. In
other words, $\alpha$ is a lower bound of $B$, but $\beta$ is not if
$\beta > \alpha$. This means that $\alpha = \inf B$.
\end{proof}

\begin{deeper}[Why this is not a triviality]
The argument runs equally well in reverse, so for ordered sets the
least-upper-bound property and the greatest-lower-bound property are
\emph{equivalent}. That is why Rudin never states a separate greatest-lower-bound
axiom for $\R$: Theorem 1.19 gives him one and Theorem 1.11 hands him the other.

Every later appearance of $\inf$ rests on this --- lower Riemann--Stieltjes sums
in 6.2, $\liminf$ in 3.16, the distance from a point to a set in the exercises.
It is a two-paragraph theorem that the rest of the book quietly leans on.
\end{deeper}

\rskip{\S\S1.12--1.18, on fields and ordered fields, are omitted from this
sample only}

\rsection{The Real Field}

\begin{signpost}[The hinge of the chapter]
Everything so far has been preparation, and everything after is consequence.
This section contains the one theorem the chapter was written to state.

\smallskip
Watch what Rudin does with it. He states 1.19, declines to prove it, and then
immediately proves 1.20 and 1.21 --- results he could have extracted from the
construction ``with very little extra effort,'' as he says. He derives them from
the least-upper-bound property instead, on purpose: he wants you to see the
property doing work before you see where it came from. Read 1.20 and 1.21 less
as facts about $\R$ than as two worked demonstrations of how a supremum is
used.
\end{signpost}

\noindent
We now state the existence theorem which is the core of this chapter.

\ritem{1.19}{Theorem}
\emph{There exists an ordered field $\R$ which has the least-upper-bound
property.}

\emph{Moreover, $\R$ contains $\Q$ as a subfield.}

The second statement means that $\Q \subset \R$ and that the operations of
addition and multiplication in $\R$, when applied to members of $\Q$, coincide
with the usual operations on rational numbers; also, the positive rational
numbers are positive elements of $\R$.\kn{``Subfield'' is a stronger claim than
``subset.'' It says the copy of $\Q$ sitting inside $\R$ has the arithmetic and
the order it had before --- so nothing you already know about rationals is
invalidated by the enlargement.}

The proof of Theorem 1.19 is rather long and a bit tedious and is therefore
presented in an Appendix to Chap.\ 1. The proof actually constructs $\R$ from
$\Q$.\kd{Worth pausing on. Everything in the next ten chapters is a consequence
of this one theorem plus the field and order axioms. If you accept 1.19, you
never need the appendix; if you doubt it, the appendix settles it by exhibiting
$\R$ as the set of Dedekind cuts of $\Q$. Either way the rest of the book is
unaffected, which is why Rudin can safely postpone it.}

The next theorem could be extracted from this construction with very little
extra effort. However, we prefer to derive it from Theorem 1.19 since this
provides a good illustration of what one can do with the least-upper-bound
property.

\ritem{1.20}{Theorem}
\begin{enumerate}[label=(\alph*), leftmargin=2.2em, itemsep=3pt, topsep=3pt]
  \item \emph{If $x \in \R$, $y \in \R$, and $x > 0$, then there is a positive
        integer $n$ such that}
        \[ nx > y. \]
  \item \emph{If $x \in \R$, $y \in \R$, and $x < y$, then there exists a
        $p \in \Q$ such that $x < p < y$.}
\end{enumerate}
Part (a) is usually referred to as the \emph{archimedean property} of $\R$. Part
(b) may be stated by saying that $\Q$ is \emph{dense} in $\R$: Between any two
real numbers there is a rational one.\kw{Part (a) is a theorem, not an axiom,
and it is not automatic. There are ordered fields in which it fails --- fields
containing an element larger than every integer. Such fields cannot have the
least-upper-bound property, and (a) is precisely the proof of that.}

\begin{proof}
(a) Let $A$ be the set of all $nx$, where $n$ runs through the positive
integers. If (a) were false, then $y$ would be an upper bound of $A$. But then
$A$ has a least upper bound in $\R$. Put $\alpha = \sup A$. Since $x > 0$,
$\alpha - x < \alpha$, and $\alpha - x$ is not an upper bound of $A$.\kn{Here is
clause (ii) of Definition 1.8 in its working form, and this is the move to
remember: to extract information from a supremum, step back from it by a
positive amount and use the fact that you are no longer an upper bound.} Hence
$\alpha - x < mx$ for some positive integer $m$. But then
$\alpha < (m+1)x \in A$, which is impossible, since $\alpha$ is an upper bound
of $A$.

(b) Since $x < y$, we have $y - x > 0$, and (a) furnishes a positive integer $n$
such that
\[ n(y-x) > 1. \]
Apply (a) again, to obtain positive integers $m_1$ and $m_2$ such that
$m_1 > nx$, $m_2 > -nx$. Then
\[ -m_2 < nx < m_1. \]
Hence there is an integer $m$ (with $-m_2 \le m \le m_1$) such that
\[ m - 1 \le nx < m. \]
If we combine these inequalities, we obtain
\[ nx < m \le 1 + nx < ny. \]
Since $n > 0$, it follows that
\[ x < \frac{m}{n} < y. \]
This proves (b), with $p = m/n$.
\end{proof}

\begin{unpack}[The three moves in part (b)]
Rudin gives the computation and not the plan. The plan is this: to fit a
rational strictly between $x$ and $y$, first \emph{stretch} the gap until it is
wider than $1$, then \emph{catch} an integer inside the stretched gap, then
\emph{shrink} back.

\smallskip
\textbf{Stretch.} Part (a) with the roles $x \leadsto y-x$ and $y \leadsto 1$
gives an $n$ with $n(y-x) > 1$. The interval $(nx, ny)$ now has length greater
than $1$.

\smallskip
\textbf{Catch.} An interval of length greater than $1$ must contain an integer.
Rudin proves this rather than asserting it, and that is what $m_1, m_2$ are for:
they trap $nx$ between two integers, so the set of integers exceeding $nx$ is
nonempty and bounded below, and therefore has a least element $m$. That gives
$m-1 \le nx < m$. The existence of a \emph{least} such integer is well-ordering
again --- the same principle used silently in 1.1. Rudin never states it as an
axiom, but this step and the archimedean property both consume it.

\smallskip
\textbf{Shrink.} From $m \le 1 + nx$ and $1 < n(y-x)$ we get
$m < nx + n(y-x) = ny$, and dividing by $n > 0$ puts $m/n$ strictly between $x$
and $y$.
\end{unpack}

\medskip
\noindent
We shall now prove the existence of $n$th roots of positive reals. This proof
will show how the difficulty pointed out in the Introduction (irrationality of
$\sqrt 2$) can be handled in $\R$.

\ritem{1.21}{Theorem}
\emph{For every real $x>0$ and every integer $n>0$ there is one and only one
positive real $y$ such that $y^n = x$.}

This number $y$ is written $\sqrt[n]{x}$ or $x^{1/n}$.

\begin{proof}
That there is at most one such $y$ is clear, since $0 < y_1 < y_2$ implies
$y_1^n < y_2^n$.

Let $E$ be the set consisting of all positive real numbers $t$ such that
$t^n < x$. If $t = x/(1+x)$ then $0 \le t < 1$. Hence $t^n \le t < x$. Thus
$t \in E$, and $E$ is not empty.\kn{Two hypotheses of the \lub{} property are
being discharged, and Rudin spends a sentence on each. This paragraph is
nonemptiness; the next is boundedness. Both are required before $\sup E$ may be
written down --- compare the warning at 1.10.}

If $t > 1+x$ then $t^n \ge t > x$, so that $t \notin E$. Thus $1+x$ is an upper
bound of $E$.

Hence Theorem 1.19 implies the existence of
\[ y = \sup E. \]

To prove that $y^n = x$ we will show that each of the inequalities $y^n < x$ and
$y^n > x$ leads to a contradiction.\kd{Note which half of Definition 1.8 each
case attacks. Assuming $y^n < x$ will contradict ``$y$ is an \emph{upper
bound}''; assuming $y^n > x$ will contradict ``$y$ is the \emph{least} upper
bound.'' Both clauses get used exactly once, which is a reliable sign that $E$
was chosen well.}

The identity $b^n - a^n = (b-a)(b^{n-1} + b^{n-2}a + \cdots + a^{n-1})$ yields
the inequality
\[ b^n - a^n < (b-a)nb^{n-1} \]
when $0 < a < b$.

Assume $y^n < x$. Choose $h$ so that $0 < h < 1$ and
\[ h < \frac{x-y^n}{n(y+1)^{n-1}}. \]
Put $a = y$, $b = y+h$. Then
\[ (y+h)^n - y^n < hn(y+h)^{n-1} < hn(y+1)^{n-1} < x - y^n. \]
Thus $(y+h)^n < x$ and $y+h \in E$. Since $y+h > y$, this contradicts the fact
that $y$ is an upper bound of $E$.

Assume $y^n > x$. Put
\[ k = \frac{y^n - x}{ny^{n-1}}. \]
Then $0 < k < y$. If $t \ge y-k$, we conclude that
\[ y^n - t^n \le y^n - (y-k)^n < kny^{n-1} = y^n - x. \]
Thus $t^n > x$ and $t \notin E$. It follows that $y - k$ is an upper bound of
$E$. But $y - k < y$, which contradicts the fact that $y$ is the least upper
bound of $E$.

Hence $y^n = x$, and the proof is complete.
\end{proof}

\begin{unpack}[Reverse-engineering the choice of $h$]
Nobody guesses that $h$. It is found by writing the goal down and solving
backwards, and Rudin has simply deleted the derivation.

\smallskip
\textbf{Goal.} Find $h>0$ with $(y+h)^n < x$; equivalently
$(y+h)^n - y^n < x - y^n$, whose right-hand side is a fixed positive number
because we assumed $y^n < x$.

\smallskip
\textbf{Bound the left side.} We cannot compute $(y+h)^n - y^n$, but the
displayed inequality with $a = y$, $b = y+h$ bounds it by $hn(y+h)^{n-1}$. So it
suffices to make $hn(y+h)^{n-1} \le x - y^n$.

\smallskip
\textbf{Get $h$ out of the exponent.} That inequality still has $h$ buried
inside $(y+h)^{n-1}$ and cannot be solved as it stands. Kill it by agreeing in
advance that $h < 1$; then $(y+h)^{n-1} < (y+1)^{n-1}$, a constant. This is the
\emph{only} reason the condition $0 < h < 1$ appears --- it is bookkeeping, not
mathematics.

\smallskip
\textbf{Solve.} It now suffices that $hn(y+1)^{n-1} < x - y^n$, that is,
$h < (x-y^n)/\bigl(n(y+1)^{n-1}\bigr)$, which is Rudin's line. Both conditions
on $h$ are upper bounds, so any $h$ below the smaller of the two works, and such
an $h$ exists.

\smallskip
The second case is the same argument run downward, and $k$ is found the same
way; there the constraint $k < y$ plays the role that $h < 1$ played here.
\end{unpack}

\begin{deeper}[What the proof is really doing]
Strip the algebra and the argument reads: \emph{$t \mapsto t^n$ is continuous,
so if $y^n$ were strictly below $x$ we could nudge $y$ upward and stay below
$x$.} The choice of $h$ is a hand-rolled continuity estimate, written out
because continuity is not available until 4.4 --- and the machinery of Chapter 4
is not available until $\R$ exists, which is what we are in the middle of
establishing.

Once Chapter 4 is in place the same theorem falls out of the intermediate value
theorem (4.23) in three lines. Rudin proves it here anyway because he needs
$n$th roots long before then, and because it is the first genuine demonstration
that the \lub{} property produces numbers $\Q$ could not.
\end{deeper}

\begin{pitfall}[Why $E$ is restricted to positive $t$]
Without positivity the uniqueness claim is simply false: for $n = 2$, $x = 4$,
both $2$ and $-2$ are roots. Rudin's ``one and only one \emph{positive} real''
is doing real work, and the restriction on $E$ is what licenses the opening line
of the proof, $0 < y_1 < y_2 \Rightarrow y_1^n < y_2^n$, which requires both
numbers to be positive.
\end{pitfall}

\begin{recap}[Chapter 1 so far, in one paragraph]
$\Q$ is unfit for analysis because a bounded set can fail to have a least upper
bound (1.1). Name that failure (1.8--1.10), assert a field free of it (1.19),
and the consequences follow at once: infima exist too (1.11), $\N$ is unbounded
and $\Q$ is dense (1.20), and every positive real has every root (1.21). Nothing
in the rest of the book uses more about $\R$ than this. When a later proof says
``by completeness,'' it means 1.19 and nothing else.
\end{recap}

\end{document}
